\chapter{Soustraction~: premiers pas}\label{ch:g1:subtraction}

Sept ballons, deux s'envolent~: combien en reste-t-il~? Enlever,
c'est \emph{soustraire}, \'ecrit avec le signe $-$. Et la soustraction
a un second visage~: elle r\'epond aussi \`a \guillemotleft{} combien
de plus~? \guillemotright{}.

\section{Enlever}

\begin{definition}[Soustraction]\label{def:g1:subtraction:def}
\emph{Soustraire}, c'est enlever et compter ce qui reste. On \'ecrit
\[
7 - 2 = 5
\]
et on lit~: \guillemotleft{} sept moins deux \'egale cinq
\guillemotright{}.
\end{definition}

\begin{omfigure}
\begin{tikzpicture}[scale=0.6]
  \foreach \i in {0,...,4} \fill[omDef] (\i*1.05,0) circle (0.3);
  \foreach \i in {5,6} {
    \fill[omDef!30] (\i*1.05,0) circle (0.3);
    \draw[omThm, thick] (\i*1.05-0.32,-0.32) -- (\i*1.05+0.32,0.32);
    \draw[omThm, thick] (\i*1.05-0.32,0.32) -- (\i*1.05+0.32,-0.32);
  }
  \node[right, font=\small] at (7.6,0)
    {$7 - 2 = 5$~: barrer deux, il en reste cinq};
\end{tikzpicture}

{\small Une soustraction en images~: barrer, c'est enlever.}
\end{omfigure}

\begin{example}[Compter \`a rebours]\label{ex:g1:subtraction:countback}
Pour une petite soustraction, compter \`a rebours sur la droite~:
$11 - 3$~: dire \guillemotleft{} onze \guillemotright{}, puis trois
pas en arri\`ere --- \guillemotleft{} dix, neuf, huit
\guillemotright{}. Donc $11 - 3 = 8$.
\end{example}

\section{Combien de plus~?}

\begin{example}[La diff\'erence]\label{ex:g1:subtraction:difference}
\guillemotleft{} Nina a $9$ billes et Tom en a $6$. Combien de billes
Nina a-t-elle de plus~? \guillemotright{} On n'enl\`eve rien, et
pourtant c'est une soustraction~: apparier les billes, et compter
celles de Nina sans partenaire~:
\[
9 - 6 = 3 .
\]
Nina en a $3$ de plus. Le r\'esultat d'une soustraction s'appelle la
\emph{diff\'erence}\index{diff\'erence}.
\end{example}

\begin{method}[Compter en montant]\label{met:g1:subtraction:countup}
Pour trouver une \omterm{ex:g1:subtraction:difference}{diff\'erence}, compter \emph{en montant} est souvent
le plus simple~: pour $12 - 9$, partir de $9$ et monter jusqu'\`a
$12$~: \guillemotleft{} dix, onze, douze \guillemotright{} --- trois
pas. Donc $12 - 9 = 3$. (Passer par $10$ pour les \'ecarts plus
grands~: $14 - 8$~: de $8$ \`a $10$ c'est $2$, de $10$ \`a $14$ c'est
$4$~: \omterm{ex:g1:subtraction:difference}{diff\'erence} $6$.)
\end{method}

\begin{omfigure}
\begin{tikzpicture}[scale=0.6]
  \draw[-Stealth] (5.6,0) -- (15.4,0);
  \foreach \i in {6,...,15} \draw (\i,-0.1) -- (\i,0.1);
  \foreach \i in {8,10,14}
    \node[below, font=\footnotesize] at (\i,-0.12) {\i};
  \draw[-Stealth, omDef, thick] (8,0.3) .. controls (9,0.9) ..
    (10,0.3);
  \node[omDef, font=\small] at (9,1.0) {$+2$};
  \draw[-Stealth, omThm, thick] (10,0.3) .. controls (12,1.2) ..
    (14,0.3);
  \node[omThm, font=\small] at (12,1.25) {$+4$};
\end{tikzpicture}

{\small $14 - 8$ en comptant en montant~: de $8$ \`a $14$ il y a
$2 + 4 = 6$ pas.}
\end{omfigure}

\section{Addition et soustraction sont jumelles}

\begin{example}[Familles de faits]\label{ex:g1:subtraction:family}
Les nombres $3$, $4$ et $7$ forment une petite famille~:
\[
3 + 4 = 7, \qquad
4 + 3 = 7, \qquad
7 - 4 = 3, \qquad
7 - 3 = 4 .
\]
Conna\^itre un fait de la famille donne les trois autres
gratuitement. La soustraction d\'efait l'addition~: apr\`es $+4$,
faire $-4$ ram\`ene au d\'epart.
\end{example}

\begin{example}[V\'erifier]\label{ex:g1:subtraction:check}
On a calcul\'e $13 - 5 = 8$~? V\'erifier avec l'addition jumelle~:
$8 + 5 = 13$. \checkmark\ Une soustraction est juste exactement quand
l'addition inverse fonctionne.
\end{example}

\begin{example}[Quelle op\'eration~?]\label{ex:g1:subtraction:choose}
\guillemotleft{} $6$ oiseaux sur le fil, $4$ s'envolent
\guillemotright{}~: enlever, $6 - 4 = 2$.
\guillemotleft{} $6$ oiseaux, puis $4$ arrivent \guillemotright{}~:
mettre ensemble, $6 + 4 = 10$.
\guillemotleft{} $6$ moineaux et $4$ rouges-gorges~: combien de
moineaux de plus~? \guillemotright{}~: \omterm{ex:g1:subtraction:difference}{diff\'erence}, $6 - 4 = 2$.
Dessiner l'histoire en cas de doute~!
\end{example}

\section{Exercices}

\begin{exercise}[$\star$]\label{exo:g1:subtraction:1}
Dessine, barre et calcule~: $6 - 2$~; \quad $8 - 5$~; \quad
$10 - 4$.
\end{exercise}

\begin{exercise}[$\star$]\label{exo:g1:subtraction:2}
Compte \`a rebours pour calculer~: $9 - 3$~; \quad $12 - 2$~; \quad
$15 - 4$.
\end{exercise}

\begin{exercise}[$\star$]\label{exo:g1:subtraction:3}
Compte en montant pour trouver la \omterm{ex:g1:subtraction:difference}{diff\'erence}~: $11 - 8$~; \quad
$13 - 9$~; \quad $16 - 12$.
\end{exercise}

\begin{exercise}[$\star$]\label{exo:g1:subtraction:4}
\'Ecris les quatre faits de la famille de $2$, $6$ et $8$.
\end{exercise}

\begin{exercise}[$\star$]\label{exo:g1:subtraction:5}
Calcule, puis v\'erifie avec l'addition jumelle~: $14 - 6$~; \quad
$17 - 9$.
\end{exercise}

\begin{exercise}[$\star$]\label{exo:g1:subtraction:6}
Compl\`ete les trous~:
\[
10 - \;?\; = 7, \qquad
\;?\; - 5 = 5, \qquad
12 - \;?\; = 12 .
\]
\end{exercise}

\begin{exercise}[$\star$]\label{exo:g1:subtraction:7}
\guillemotleft{} Il y a $15$ cerises dans le bol. Les enfants en
mangent $8$. Combien en reste-t-il~? \guillemotright{} \'Ecris la
soustraction et la phrase r\'eponse.
\end{exercise}

\begin{exercise}[$\star$]\label{exo:g1:subtraction:8}
Pour chaque histoire, choisis $+$ ou $-$, puis calcule~:
\guillemotleft{} $9$ canards sur l'\'etang, $5$ s'\'eloignent
\guillemotright{}~; \guillemotleft{} $9$ canards, $5$ oies
arrivent \guillemotright{}~; \guillemotleft{} Ali a $9$~ans, sa
s{\oe}ur en a $5$~: de combien d'ann\'ees Ali est-il plus \^ag\'e~?
\guillemotright{}.
\end{exercise}

\begin{exercise}[$\star\star$]\label{exo:g1:subtraction:9}
Maman cane avait $12$ {\oe}ufs. Certains ont \'eclos~; maintenant
$12 - ?$ {\oe}ufs restent et elle compte $7$ {\oe}ufs non \'eclos.
Combien d'{\oe}ufs ont \'eclos~? (Quelle famille de faits est-ce~?)
\end{exercise}
